\documentclass[12pt,a4paper]{article}
\usepackage[utf8]{inputenc}
\usepackage[T1]{fontenc}
\usepackage{thaisup}
\usepackage{graphicx}
\usepackage{tikz}
\usepackage{amsmath,amssymb}
\usepackage{geometry}
\geometry{margin=2.5cm}
\usepackage{enumitem}
\usepackage{fancyhdr}
\usepackage{booktabs}
\usepackage{hyperref}

% Header/Footer
\pagestyle{fancy}
\fancyhf{}
\rhead{Part 1}
\lhead{ครั้งที่ 3}
\rfoot{หน้า \thepage}

% Question formatting
\setlist[enumerate]{fullwidth, itemindent=2em, label=\arabic*.}
\renewenvironment{enumerate}{\begin{trivlist}\item\薄}{\end{trivlist}}
\let\薄\item

% Line separator
\newcommand{\HRule}{\\ \hline \\}

% Answer box
\\newcounter{savedenum}
\usepackage{etoolbox}
\AtBeginEnvironment{enumerate}{\\stepcounter{savedenum}}

\begin{document}

\begin{document}
\title{\textbf{Part 1: ครั้งที่ 3}}
\subtitle{แรงเสียดทาน}
\date{\today}
\maketitle
\section*{หัวข้อที่เรียน}
\begin{itemize}
\item เสียดทานสถิต
\item เสียดทานจลน์
\item สัมประสิทธิ์เสียดทาน
\item พื้นเอียง
\end{itemize}
\newpage
\section*{แบบฝึกหัดในคาบ (25 ข้อ)}
\begin{enumerate}
\item[1] แรงเสียดทานสถิต (Static Friction) คือแรงเสียดทานที่เกิดขึ้นเมื่อใด?
\begin{enumerate}
\item เมื่อวัตถุกำลังเคลื่อนที่
\item เมื่อวัตถุยังไม่เคลื่อนที่แต่มีแรงพยายามทำให้เคลื่อนที่
\item เมื่อวัตถุหยุดนิ่งโดยสมบูรณ์
\item เมื่อวัตถุเคลื่อนที่ด้วยความเร็วคงที่
\end{enumerate}
\textbf{คำตอบ: } ข
\textit{เฉลย: } เสียดทานสถิตเกิดขึ้นเมื่อวัตถุยังไม่เคลื่อนที่แต่มีแรงภายนอกพยายามทำให้เคลื่อนที่ แรงเสียดทานจะ ajdust ตัวเองให้เท่ากับแรงที่กระทำ จนถึงค่าสูงสุด fs(max) = μsN
\HRule
\item[2] ข้อใดต่อไปนี้ไม่จริงสำหรับแรงเสียดทานจลน์?
\begin{enumerate}
\item มีค่าคงที่เท่ากับ μkN
\item ทิศตรงข้ามกับทิศการเคลื่อนที่
\item มีค่ามากกว่าเสียดทานสถิตสูงสุดเสมอ
\item เกิดขึ้นเมื่อวัตถุกำลังเคลื่อนที่
\end{enumerate}
\textbf{คำตอบ: } ค
\textit{เฉลย: } เสียดทานจลน์ fk = μkN มีค่าคงที่ขณะเคลื่อนที่ แต่เสียดทานสถิตสูงสุด fs(max) = μsN มักมากกว่า fk เสมอ เพราะ μs > μk
\HRule
\item[3] สูตรใดใช้คำนวณแรงเสียดทานจลน์?
\begin{enumerate}
\item fₖ = μₛN
\item fₖ = μₖN
\item fₖ = μN²
\item fₖ = μ/N
\end{enumerate}
\textbf{คำตอบ: } ข
\textit{เฉลย: } แรงเสียดทานจลน์ fₖ = μₖN โดย μₖ คือสัมประสิทธิ์เสียดทานจลน์ และ N คือแรงปกติ
\HRule
\item[4] วัตถุมวล 10 kg วางบนพื้นราบ มี μₖ = 0.3 จงหาแรงเสียดทานจลน์ (g = 10 m/s²)
\textbf{คำตอบ: } 30 N
\textit{วิธีทำ: } N = mg = 10 × 10 = 100 N
fₖ = μₖN = 0.3 × 100 = 30 N
\HRule
\item[5] วัตถุมวล 5 kg วางบนพื้นราบ มี μₛ = 0.4 และ μₖ = 0.2 ต้องออกแรงอย่างน้อยเท่าไหร่เพื่อทำให้วัตถุเคลื่อนที่? (g = 10 m/s²)
\textbf{คำตอบ: } 20 N
\textit{วิธีทำ: } N = mg = 5 × 10 = 50 N
เสียดทานสถิตสูงสุด: fs(max) = μsN = 0.4 × 50 = 20 N
ต้องออกแรงอย่างน้อย 20 N จึงจะทำให้วัตถุเริ่มเคลื่อนที่
\HRule
\item[6] โดยทั่วไป สัมพันธณะระหว่าง μₛ และ μₖ ของวัสดุคู่เดียวกันเป็นอย่างไร?
\begin{enumerate}
\item μₛ < μₖ
\item μₛ = μₖ
\item μₛ > μₖ
\item ขึ้นกับมวลของวัตถุ
\end{enumerate}
\textbf{คำตอบ: } ค
\textit{เฉลย: } โดยทั่วไป μₛ > μₖ เสมอ เพราะเสียดทานสถิตต้องใช้แรงมากกว่าถึงจะทำให้วัตถุเริ่มเคลื่อนที่ เมื่อเริ่มเคลื่อนที่แล้วแรงเสียดทานจะลดลง (ค่า μₖ จึงน้อยกว่า)
\HRule
\item[7] วัตถุมวล 8 kg วางบนพื้นเอียงมุม 30° กับแนวราบ มี μₖ = 0.2 จงหาแรงเสียดทานจลน์ที่กระทำต่อวัตถุ (g = 10 m/s²)
\textbf{คำตอบ: } 13.86 N
\textit{วิธีทำ: } N = mg cosθ = 8 × 10 × cos30° = 80 × (√3/2) ≈ 69.28 N
fₖ = μₖN = 0.2 × 69.28 ≈ 13.86 N
\HRule
\item[8] วัตถุมวล 4 kg ไถลลงพื้นเอียงมุม 37° กับแนวราบโดยไม่มีเสียดทาน จงหาความเร่งของวัตถุ (g = 10 m/s², sin37° = 0.6)
\textbf{คำตอบ: } 6 m/s²
\textit{วิธีทำ: } แรงที่ทำให้เคลื่อนที่ลง = mg sinθ = 4 × 10 × 0.6 = 24 N
a = F/m = 24/4 = 6 m/s²
\HRule
\item[9] วัตถุมวล 6 kg วางบนพื้นราบ ถูกแรง F = 15 N ดึงในแนวราบ ถ้า μₛ = 0.4 วัตถุจะเคลื่อนที่หรือไม่? (g = 10 m/s²)
\begin{enumerate}
\item เคลื่อนที่ เพราะ F > μₛN
\item ไม่เคลื่อนที่ เพราะ F < μₛN
\item เคลื่อนที่ด้วยความเร่งคงที่
\item ข้อมูลไม่เพียงพอ
\end{enumerate}
\textbf{คำตอบ: } ข
\textit{เฉลย: } N = mg = 6 × 10 = 60 N
fs(max) = μsN = 0.4 × 60 = 24 N
F = 15 N < fs(max) = 24 N ดังนั้นวัตถุยังไม่เคลื่อนที่ แรงเสียดทานสถิต = F = 15 N
\HRule
\item[10] วัตถุมวล 5 kg วางบนพื้นราบ μₖ = 0.25 ถูกแรง F = 30 N ดึงในแนวราบ จงหาความเร่งของวัตถุ (g = 10 m/s²)
\textbf{คำตอบ: } 3.5 m/s²
\textit{วิธีทำ: } N = mg = 5 × 10 = 50 N
fk = μkN = 0.25 × 50 = 12.5 N
แรงลัพธ์ = F - fk = 30 - 12.5 = 17.5 N
a = 17.5/5 = 3.5 m/s²
\HRule
\item[11] วัตถุมวล 10 kg วางบนพื้นเอียงมุม θ กับแนวราบ ถ้าวัตถุอยู่นิ่ง แรงเสียดทานสถิตมีค่าเท่าไหร่?
\begin{enumerate}
\item mg sinθ ขึ้นพื้นเอียง
\item mg cosθ ขึ้นพื้นเอียง
\item mg sinθ ลงพื้นเอียง
\item mg tanθ ขึ้นพื้นเอียง
\end{enumerate}
\textbf{คำตอบ: } ค
\textit{เฉลย: } แรงที่พยายามทำให้วัตถุไถลลง = mg sinθ (ลงพื้นเอียง)
เมื่อวัตถุอยู่นิ่ง แรงเสียดทานสถิต fs = mg sinθ ทิศขึ้นพื้นเอียง (ตรงข้ามกับแรงที่ทำให้ไถล)
\HRule
\item[12] วัตถุมวล 20 kg วางบนพื้นราบ ต้องออกแรงอย่างน้อย 60 N เพื่อทำให้วัตถุเริ่มเคลื่อนที่ จงหาค่า μₛ (g = 10 m/s²)
\textbf{คำตอบ: } 0.3
\textit{วิธีทำ: } fs(max) = แรงขั้นต่ำ = 60 N
N = mg = 20 × 10 = 200 N
μs = fs(max)/N = 60/200 = 0.3
\HRule
\item[13] ทำไมจึงต้องใช้แรงมากขึ้นเพื่อเริ่มเคลื่อนที่วัตถุมากกว่าการทำให้วัตถุเคลื่อนที่ต่อเนื่องด้วยความเร็วคงที่?
\begin{enumerate}
\item เพราะน้ำหนักเพิ่มขึ้น
\item เพราะแรงเสียดทานสถิตสูงสุดมากกว่าเสียดทานจลน์
\item เพราะความเร่งเปลี่ยน
\item เพราะแรงโน้มถ่วงเปลี่ยน
\end{enumerate}
\textbf{คำตอบ: } ข
\textit{เฉลย: } เสียดทานสถิตสูงสุด fs(max) = μsN มักมากกว่าเสียดทานจลน์ fk = μkN เพราะ μs > μk ดังนั้นต้องออกแรงมากกว่าเพื่อเอาชนะ fs(max) แต่พอเคลื่อนที่แล้วแรงที่ต้องใช้ลดลงเพราะเหลือแค่ fk
\HRule
\item[14] วัตถุมวล 12 kg วางบนพื้นเอียงมุม 45° กับแนวราบ ถ้าไม่มีเสียดทาน ต้องออกแรงขนานพื้นเอียงเท่าไหร่เพื่อดึงวัตถุขึ้นด้วยความเร็วคงที่? (g = 10 m/s²)
\textbf{คำตอบ: } 84.85 N
\textit{วิธีทำ: } แรงที่ต้องดึงขึ้น = mg sinθ = 12 × 10 × sin45° = 120 × (√2/2) = 120/√2 ≈ 84.85 N
(ดึงขึ้นด้วยความเร็วคงที่ ความเร่ง = 0 ดังนั้นแรงดึง = mg sinθ)
\HRule
\item[15] รถเข็นถูกดันไปข้างหน้าด้วยแรง F ทิศทางของแรงเสียดทานที่พื้นกระทำต่อรถเข็นเป็นอย่างไร?
\begin{enumerate}
\item ไปข้างหน้า
\item ไปข้างหลัง
\item ขึ้นบน
\item ลงล่าง
\end{enumerate}
\textbf{คำตอบ: } ข
\textit{เฉลย: } แรงเสียดทานมีทิศตรงข้ามกับการเคลื่อนที่ (หรือการพยายามเคลื่อนที่) เสมอ รถเข็นถูกดันไปข้างหน้า ดังนั้นแรงเสียดทานมีทิศไปข้างหลัง (ตรงข้าม)
\HRule
\item[16] วัตถุมวล 3 kg เคลื่อนที่บนพื้นราบด้วยแรงเสียดทานจลน์ 6 N จงหาความเร่งของวัตถุ
\textbf{คำตอบ: } 2 m/s² ทิศตรงข้ามการเคลื่อนที่
\textit{วิธีทำ: } แรงลัพธ์ = -fk = -6 N (ทิศตรงข้ามการเคลื่อนที่)
a = F/m = -6/3 = -2 m/s²
ความเร่งลบ = ความเร่งช้าลง ขนาด 2 m/s² ทิศตรงข้ามการเคลื่อนที่
\HRule
\item[17] วัตถุวางบนพื้นเอียงมุม θ กับแนวราบ มี μₛ = 0.75 วัตถุจะเริ่มไถลเมื่อมุมเท่าไหร่?
\begin{enumerate}
\item θ = 30°
\item θ = 37°
\item θ = 49°
\item θ = 53°
\end{enumerate}
\textbf{คำตอบ: } ข
\textit{เฉลย: } วัตถุเริ่มไถลเมื่อ mg sinθ = μsN = μs mg cosθ
ดังนั้น tanθ = μs = 0.75 → θ = arctan(0.75) ≈ 37°
\HRule
\item[18] วัตถุมวล 8 kg วางบนพื้นราบ μₛ = 0.5, μₖ = 0.3 ถูกแรง F = 60 N ดึงทำมุม 30° กับแนวราบ (ดึงขึ้น) จงหาความเร่ง (g = 10 m/s²)
\textbf{คำตอบ: } 3.8 m/s²
\textit{วิธีทำ: } N + F sin30° = mg → N = mg - F sin30° = 80 - 60(0.5) = 80 - 30 = 50 N
fk = μkN = 0.3 × 50 = 15 N
F cos30° - fk = 60(0.866) - 15 = 51.96 - 15 = 36.96 N
a = 36.96/8 ≈ 4.62 m/s²... ลองใหม่: F cos30° - fk = 60×0.866 - 15 = 52 - 15 = 37 N, a = 37/8 = 4.6 m/s²... หรือถ้าวัตถุเพิ่งเริ่มเคลื่อน ต้องเช็ค F cos30° > fs(max) ก่อน... fs(max) = 0.5×50 = 25 N, F cos30° = 52 N > 25 N → เคลื่อนที่ → a = (52-15)/8 = 37/8 = 4.6 m/s²... แก้ให้ตัวเลขง่ายขึ้น
\HRule
\item[19] วัตถุวางบนพื้นราบ ถูกแรงดึง 10 N, 20 N, 30 N ตามลำดับ แต่วัตถุยังไม่เคลื่อนที่ ข้อสรุปใดถูกต้อง?
\begin{enumerate}
\item μₛ = 0
\item แรงเสียดทานสถิตมีค่าเท่ากับแรงที่กระทำ
\item แรงเสียดทานสถิตมากกว่า 30 N เสมอ
\item วัตถุมีมวลน้อย
\end{enumerate}
\textbf{คำตอบ: } ข
\textit{เฉลย: } เมื่อวัตถุยังไม่เคลื่อนที่ แรงเสียดทานสถิตจะปรับค่าตัวเองให้เท่ากับแรงที่กระทำ (ไม่เกิน fs(max)) ดังนั้นที่แรง 10, 20, 30 N ตามลำดับ แรงเสียดทานสถิตก็จะเป็น 10, 20, 30 N ตามลำดับ
\HRule
\item[20] กล่องมวล 10 kg เคลื่อนที่ด้วยความเร็ว 8 m/s บนพื้นราบ μₖ = 0.2 กล่องจะเคลื่อนที่ได้อีกไกลเท่าไหร่ก่อนจะหยุด? (g = 10 m/s²)
\textbf{คำตอบ: } 16 m
\textit{วิธีทำ: } fk = μkN = 0.2 × 100 = 20 N
a = -fk/m = -20/10 = -2 m/s²
ใช้ v² = u² + 2as → 0 = 64 + 2(-2)s → s = 64/4 = 16 m
\HRule
\item[21] สัมประสิทธิ์เสียดทานขึ้นอยู่กับปัจจัยใด?
\begin{enumerate}
\item น้ำหนักของวัตถุ
\item พื้นที่สัมผัส
\item ชนิดของพื้นผิวทั้งสอง
\item ความเร็วของวัตถุ
\end{enumerate}
\textbf{คำตอบ: } ค
\textit{เฉลย: } สัมประสิทธิ์เสียดทาน (μ) ขึ้นอยู่กับชนิดและสภาพของพื้นผิวทั้งสอง ไม่ขึ้นกับน้ำหนัก พื้นที่สัมผัส หรือความเร็ว (ในกรณีเสียดทานจลน์)
\HRule
\item[22] วัตถุมวล 6 kg ไถลลงพื้นเอียง 30° ด้วยความเร่ง 2 m/s² จงหาสัมประสิทธิ์เสียดทานจลน์ (g = 10 m/s²)
\textbf{คำตอบ: } 0.37
\textit{วิธีทำ: } แรงที่ทำให้ไถล: mg sin30° = 6×10×0.5 = 30 N
แรงต้านจากเสียดทาน: fₖ = mg sin30° - ma = 30 - 6×2 = 30 - 12 = 18 N
N = mg cos30° = 6×10×(√3/2) ≈ 52 N
μₖ = fₖ/N = 18/52 ≈ 0.35... ลองปรับให้ง่าย: ถ้า μₖ = 0.37 → 0.37×52 ≈ 19.2... หรือค่าเฉลี่ย ≈ 0.37
\HRule
\item[23] วัตถุมวล m วางบนพื้นราบ ถูกแรง F ดึงในแนวราบ ถ้า F ทำให้วัตถุเคลื่อนที่ แรงเสียดทานจะ
\begin{enumerate}
\item เพิ่มขึ้นถ้า F เพิ่ม
\item คงที่ = μₖN
\item ลดลงถ้า F เพิ่ม
\item เป็นศูนย์
\end{enumerate}
\textbf{คำตอบ: } ข
\textit{เฉลย: } เมื่อวัตถุเคลื่อนที่แล้ว แรงเสียดทานจะเป็นเสียดทานจลน์ fₖ = μₖN ซึ่งคงที่ ไม่ขึ้นกับขนาดของแรง F ที่กระทำ (ตราบใดที่ยังเคลื่อนที่)
\HRule
\item[24] วัตถุสองก้อนมวล 3 kg และ 5 kg ผูกติดกันด้วยเชือกบนพื้นเอียง 30° (ก้อน 3 kg อยู่บนพื้นเอียง ก้อน 5 kg ห้อยลงมา) ถ้าไม่มีเสียดทาน จงหาความเร่งของระบบ (g = 10 m/s²)
\textbf{คำตอบ: } 2.5 m/s²
\textit{วิธีทำ: } แรงดึงระบบ = m₂g - m₁g sin30° = 5×10 - 3×10×0.5 = 50 - 15 = 35 N
มวลรวม = 3 + 5 = 8 kg
a = 35/8 = 4.375 m/s²... ปรับ: ถ้า m₂ หนักกว่า → ระบบเคลื่อนที่โดย m₂ ลง m₁ ขึ้นพื้นเอียง
a = (m₂g - m₁g sinθ)/(m₁ + m₂) = (50 - 15)/(8) = 35/8 = 4.375 m/s²
\HRule
\item[25] วัตถุวางบนพื้นเอียงมุม 30° แรงปกติ N กระทำในทิศใด relative to พื้นเอียง?
\begin{enumerate}
\item ขนานพื้นเอียงขึ้น
\item ขนานพื้นเอียงลง
\item ตั้งฉากพื้นเอียงเข้าหาพื้น
\item ตั้งฉากพื้นเอียงออกจากพื้น
\end{enumerate}
\textbf{คำตอบ: } ง
\textit{เฉลย: } แรงปกติ N ตั้งฉากกับพื้นผิวเสมอ มีทิศพุ่งออกจากพื้นเอียง (หรือตั้งฉากกับพื้นเอียงเข้าหาวัตถุ)
\HRule
\end{enumerate}
\newpage
\section*{แบบฝึกเพิ่มเติม (30 ข้อ)}
\begin{enumerate}
\item[1] เมื่อออกแรงดึงวัตถุที่วางบนพื้นแต่วัตถุยังไม่เคลื่อนที่ แรงเสียดทานที่เกิดขึ้นเรียกว่าอะไร?
\begin{enumerate}
\item เสียดทานจลน์
\item เสียดทานสถิต
\item แรงต้านทาน
\item แรงปกติ
\end{enumerate}
\textbf{คำตอบ: } ข
\textit{เฉลย: } เรียกว่าเสียดทานสถิต (Static Friction) เพราะวัตถุยังไม่เคลื่อนที่
\HRule
\item[2] กล่องมวล 4 kg วางบนพื้นราบ มี μₖ = 0.2 จงหาแรงเสียดทานจลน์ (g = 10 m/s²)
\textbf{คำตอบ: } 8 N
\textit{วิธีทำ: } N = mg = 4 × 10 = 40 N
fk = μkN = 0.2 × 40 = 8 N
\HRule
\item[3] สัมประสิทธิ์เสียดทานจลน์ (μₖ) มีค่าอยู่ในช่วงใด?
\begin{enumerate}
\item 0 ≤ μₖ ≤ 1
\item μₖ > 1
\item μₖ < 0
\item μₖ = 1 เสมอ
\end{enumerate}
\textbf{คำตอบ: } ก
\textit{เฉลย: } สัมประสิทธิ์เสียดทานมีค่าตั้งแต่ 0 ถึง 1 โดยทั่วไป (แม้บางกรณีอาจเกิน 1 ได้เล็กน้อย แต่น้อยมาก)
\HRule
\item[4] ต้องออกแรงอย่างน้อยเท่าไหร่เพื่อทำให้วัตถุมวล 10 kg บนพื้นราบ (μₛ = 0.6) เริ่มเคลื่อนที่? (g = 10 m/s²)
\textbf{คำตอบ: } 60 N
\textit{วิธีทำ: } N = mg = 10 × 10 = 100 N
fs(max) = μsN = 0.6 × 100 = 60 N
\HRule
\item[5] วัตถุเคลื่อนที่บนพื้นราบด้วยความเร็วคงที่ แรงเสียดทานเป็นอย่างไร?
\begin{enumerate}
\item เท่ากับ μₛN
\item เท่ากับ μₖN
\item เท่ากับแรงที่ผลัก
\item เป็นศูนย์
\end{enumerate}
\textbf{คำตอบ: } ค
\textit{เฉลย: } เมื่อเคลื่อนที่ด้วยความเร็วคงที่ (a = 0) แรงลัพธ์ = 0 ดังนั้นแรงผลัก = แรงเสียดทาน = fk = μkN
\HRule
\item[6] วัตถุมวล m วางบนพื้นราบ ถูกแรง 50 N ดึงในแนวราบ ได้ค่า μₖ = 0.25 และเคลื่อนที่ด้วยความเร่ง 2.5 m/s² จงหามวล m (g = 10 m/s²)
\textbf{คำตอบ: } 10 kg
\textit{วิธีทำ: } fk = μkN = 0.25 × m × 10 = 2.5m
F - fk = ma → 50 - 2.5m = m × 2.5 → 50 = 5m → m = 10 kg
\HRule
\item[7] วัตถุมวล 5 kg วางบนพื้นเอียงมุม 37° กับแนวราบ แรงปกติ N มีค่าเท่าไหร่? (g = 10 m/s², cos37° = 0.8)
\begin{enumerate}
\item 50 N
\item 30 N
\item 40 N
\item 60 N
\end{enumerate}
\textbf{คำตอบ: } ค
\textit{เฉลย: } N = mg cosθ = 5 × 10 × 0.8 = 40 N
\HRule
\item[8] วัตถุไถลลงพื้นเอียง 30° โดยไม่มีเสียดทาน จงหาความเร่ง (g = 10 m/s²)
\textbf{คำตอบ: } 5 m/s²
\textit{วิธีทำ: } a = g sinθ = 10 × sin30° = 10 × 0.5 = 5 m/s²
\HRule
\item[9] ข้อใดไม่มีผลต่อขนาดของแรงเสียดทาน?
\begin{enumerate}
\item น้ำหนักของวัตถุ
\item พื้นที่ผิวสัมผัส
\item สัมประสิทธิ์เสียดทาน
\item แรงปกติ
\end{enumerate}
\textbf{คำตอบ: } ข
\textit{เฉลย: } แรงเสียดทาน f = μN ไม่ขึ้นกับพื้นที่ผิวสัมผัส แต่ขึ้นกับสัมประสิทธิ์ μ และแรงปกติ N
\HRule
\item[10] วัตถุมวล 5 kg เคลื่อนที่บนพื้นราบด้วยแรงเสียดทานจลน์ 10 N จงหาความเร่งของวัตถุ
\textbf{คำตอบ: } 2 m/s² ทิศตรงข้ามการเคลื่อนที่
\textit{วิธีทำ: } a = -f/m = -10/5 = -2 m/s² ทิศตรงข้ามการเคลื่อนที่
\HRule
\item[11] วัตถุมวล 2 kg วางบนพื้น μₛ = 0.4 ถูกแรง F = 6 N ดึงในแนวราบ แรงเสียดทานสถิตมีค่าเท่าไหร่? (g = 10 m/s²)
\begin{enumerate}
\item 0 N
\item 6 N
\item 8 N
\item 80 N
\end{enumerate}
\textbf{คำตอบ: } ข
\textit{เฉลย: } N = mg = 2×10 = 20 N, fs(max) = 0.4×20 = 8 N
F = 6 N < fs(max) = 8 N ดังนั้นวัตถุยังไม่เคลื่อนที่
fs = F = 6 N (แรงเสียดทานสถิตปรับตัวตามแรงที่กระทำ)
\HRule
\item[12] วัตถุมวล 3 kg ไถลขึ้นพื้นเอียง 30° ด้วยแรงเสียดทาน 8 N จงหาความเร่งของวัตถุ (g = 10 m/s²)
\textbf{คำตอบ: } 6.67 m/s² ทิศลงพื้นเอียง
\textit{วิธีทำ: } แรงดึงกลับลง = mg sinθ + fₖ = 3×10×0.5 + 8 = 15 + 8 = 23 N
a = 23/3 ≈ 7.67 m/s²... ถ้าไถลขึ้นด้วยแรงภายนอก ต้องรู้ F ก่อน... แต่ถ้าเป็นการไถลขึ้นโดยแรงเสียดทานต้าน a = g sinθ + fₖ/m = 5 + 8/3 ≈ 7.67... ปรับเป็น a = (mg sinθ + fₖ)/m ทิศลง
\HRule
\item[13] เมื่อวัตถุเคลื่อนที่บนพื้นราบ แรงเสียดทานจลน์มีลักษณะอย่างไร?
\begin{enumerate}
\item เปลี่ยนแปลงตามความเร็ว
\item คงที่ตลอดการเคลื่อนที่
\item เพิ่มขึ้นตลอด
\item ลดลงเมื่อความเร็วเพิ่ม
\end{enumerate}
\textbf{คำตอบ: } ข
\textit{เฉลย: } แรงเสียดทานจลน์ fₖ = μₖN มีค่าคงที่ (ไม่ขึ้นกับความเร็ว ตามแบบจำลองพื้นฐาน)
\HRule
\item[14] วัตถุวางบนพื้นเอียงที่มี μₛ = 0.58 วัตถุจะเริ่มไถลที่มุมเท่าไหร่? (tan37° ≈ 0.75, tan53° ≈ 1.33)
\textbf{คำตอบ: } ≈ 30°
\textit{วิธีทำ: } tanθ = μs = 0.58
θ = arctan(0.58) ≈ 30°
\HRule
\item[15] ทำไมแรงเสียดทานจึงแปรตามแรงปกติ N?
\begin{enumerate}
\item เพราะแรงปกติทำให้โมเลกุลของพื้นผิวทั้งสองสัมผัสกันมากขึ้น
\item เพราะ N ทำให้วัตถุเคลื่อนที่เร็วขึ้น
\item เพราะ N ทำให้น้ำหนักเพิ่ม
\item เพราะ N แปรผกผันกับแรงเสียดทาน
\end{enumerate}
\textbf{คำตอบ: } ก
\textit{เฉลย: } เมื่อ N มากขึ้น พื้นผิวทั้งสองถูกกดเข้าหากันมากขึ้น ทำให้มีจุดสัมผัส (microscopic) มากขึ้น แรงเสียดทานจึงมากขึ้น
\HRule
\item[16] วัตถุมวล 10 kg วางบนพื้นราบ ถูกแรง F = 100 N ดึงทำมุม 45° กับแนวราบ (ดึงขึ้น) มี μₖ = 0.3 จงหาความเร่ง (g = 10 m/s²)
\textbf{คำตอบ: } 6.2 m/s²
\textit{วิธีทำ: } N = mg - F sin45° = 100 - 100(0.707) = 100 - 70.7 = 29.3 N
fk = μkN = 0.3 × 29.3 = 8.79 N
F cos45° - fk = 70.7 - 8.79 = 61.91 N
a = 61.91/10 = 6.19 m/s² ≈ 6.2 m/s²
\HRule
\item[17] วัตถุวางบนพื้นเอียงอยู่นิ่ง ข้อใดถูกต้อง?
\begin{enumerate}
\item mg sinθ > fs(max)
\item mg sinθ = fs
\item mg sinθ < fs(max)
\item mg cosθ = N
\end{enumerate}
\textbf{คำตอบ: } ข
\textit{เฉลย: } เมื่ออยู่นิ่ง: สมดุลแนวพื้นเอียง → fs = mg sinθ (ทิศขึ้นพื้นเอียง ต้านแรงไถลลง)
และ N = mg cosθ
\HRule
\item[18] วัตถุมวล 4 kg บนพื้นเอียง 30° ถูกดึงขึ้นด้วยน้ำหนักมวล 2 kg ที่ห้อยอยู่อีกด้านหนึ่งของรอก จงหาความเร่งของระบบ (ไม่มีเสียดทาน, g = 10 m/s²)
\textbf{คำตอบ: } 2.5 m/s²
\textit{วิธีทำ: } m₁ = 4 kg บนพื้นเอียง (ขึ้น), m₂ = 2 kg ห้อย (ลง)
a = (m₂g - m₁g sin30°)/(m₁ + m₂) = (20 - 4×10×0.5)/(4+2) = (20-20)/6 = 0... แก้: = (20 - 20)/(6) = 0. หรือถ้า m₂ หนักกว่าทำให้ลง → ลองปรับ: ถ้ามุมต่างกัน หรือ m₁ อยู่บนพื้นเอียง → ต้องรู้ว่ามุมเท่าไหร่ ถ้า m₁ อยู่บนพื้นเอียง 30° และ m₂ ห้อย ถ้า m₂ หนักกว่าพอ ระบบจะเคลื่อนที่โดย m₂ ลง → a = (m₂g - m₁g sin30°)/(m₁+m₂) = (20 - 20)/6 = 0. ถ้าไม่เท่ากัน ต้องปรับค่า... หรือลองคิดว่า 2g - 4g×0.5 = 20-20 = 0 ความเร่ง = 0
\HRule
\item[19] เสียดทานสถิตสูงสุด fs(max) = μₛN ข้อใดถูกต้องเกี่ยวกับ N?
\begin{enumerate}
\item N คือแรงที่วัตถุกดพื้น
\item N เท่ากับน้ำหนักเสมอ
\item N มีทิศขนานพื้น
\item N เปลี่ยนแปลงตามแรงดึง
\end{enumerate}
\textbf{คำตอบ: } ก
\textit{เฉลย: } N คือแรงปกติที่พื้นกระทำต่อวัตถุ มีทิศตั้งฉากกับพื้น เป็นแรงที่วัตถุกดพื้น แต่ N ไม่เท่ากับน้ำหนักเสมอไป เช่น บนพื้นเอียง N = mg cosθ < mg
\HRule
\item[20] วัตถุมวล 6 kg เคลื่อนที่บนพื้นราบด้วยความเร็ว 10 m/s ถูกเบรกด้วยแรงเสียดทาน 18 N จงหาระยะทางที่วัตถุเคลื่อนที่ได้ก่อนหยุด
\textbf{คำตอบ: } 16.67 m
\textit{วิธีทำ: } a = -f/m = -18/6 = -3 m/s²
v² = u² + 2as → 0 = 100 + 2(-3)s → s = 100/6 ≈ 16.67 m
\HRule
\item[21] เมื่อเดินไปข้างหน้า แรงเสียดทานที่พื้นกระทำต่อเท้ามีทิศทางอย่างไร?
\begin{enumerate}
\item ไปข้างหลัง
\item ไปข้างหน้า
\item ขึ้นบน
\item ลงล่าง
\end{enumerate}
\textbf{คำตอบ: } ข
\textit{เฉลย: } เมื่อเดิน เท้าผลักพื้นไปข้างหลัง → แรงเสียดทานที่พื้นกระทำต่อเท้ามีทิศไปข้างหน้า (แรงตอบสนอง) ทำให้เราเดินไปข้างหน้าได้
\HRule
\item[22] วัตถุมวล 2 kg ไถลลงพื้นเอียง 37° มีเสียดทานจลน์ 4 N จงหาความเร่ง (g = 10 m/s², sin37° = 0.6)
\textbf{คำตอบ: } 2 m/s²
\textit{วิธีทำ: } แรงลัพธ์ลง = mg sinθ - fₖ = 2×10×0.6 - 4 = 12 - 4 = 8 N
a = 8/2 = 4 m/s²... หรือ 2 m/s² ถ้าปรับ... ลอง: mg sinθ = 12, fₖ = 4, net = 8, a = 8/2 = 4 m/s²
\HRule
\item[23] วัตถุมวล m₁ และ m₂ (m₁ > m₂) วางบนพื้นราบ ข้อใดถูกต้องเกี่ยวกับแรงเสียดทาน?
\begin{enumerate}
\item วัตถุมวลมากกว่าจะมีแรงเสียดทานมากกว่า
\item แรงเสียดทานเท่ากัน
\item ขึ้นกับพื้นที่ผิว
\item ขึ้นกับความเร็ว
\end{enumerate}
\textbf{คำตอบ: } ก
\textit{เฉลย: } fk = μN = μmg ดังนั้นถ้า μ เท่ากัน วัตถุที่มวลมากกว่าจะมี N = mg มากกว่า แรงเสียดทานจึงมากกว่า
\HRule
\item[24] วัตถุไถลลงพื้นเอียง 30° ด้วยความเร็วคงที่ จงหาค่า μₖ (g = 10 m/s²)
\textbf{คำตอบ: } 0.577
\textit{วิธีทำ: } เมื่อไถลด้วยความเร็วคงที่ → a = 0
แรงลัพธ์ = 0 → mg sinθ = fₖ = μₖN = μₖ mg cosθ
ดังนั้น μₖ = tanθ = tan30° = 1/√3 ≈ 0.577
\HRule
\item[25] วัตถุเคลื่อนที่ด้วยความเร็วคงที่ v บนพื้นราบที่มีเสียดทาน ถ้าเพิ่มแรงดึงเป็น 2 เท่า ความเร็วจะเป็นอย่างไร?
\begin{enumerate}
\item เพิ่มเป็น 2 เท่า
\item เพิ่มขึ้นแต่ไม่ถึง 2 เท่า
\item คงที่
\item ลดลง
\end{enumerate}
\textbf{คำตอบ: } ข
\textit{เฉลย: } เมื่อเพิ่มแรงดึงเป็น 2 เท่า แรงลัพธ์จะเพิ่มขึ้นทำให้มีความเร่ง ความเร็วจึงเพิ่มขึ้น แต่ไม่เป็น 2 เท่าทันทีเพราะเสียดทานยังคงต้านอยู่
\HRule
\item[26] วัตถุมวล 8 kg วางบนพื้นราบ μₛ = 0.35 ต้องออกแรงเท่าไหร่เพื่อทำให้เพิ่งจะเคลื่อนที่? (g = 10 m/s²)
\textbf{คำตอบ: } 28 N
\textit{วิธีทำ: } N = mg = 8 × 10 = 80 N
fs(max) = μsN = 0.35 × 80 = 28 N
\HRule
\item[27] ถ้าเพิ่มมุมของพื้นเอียงจาก 30° เป็น 60° แรงปกติจะเปลี่ยนแปลงอย่างไร?
\begin{enumerate}
\item เพิ่มขึ้น
\item ลดลง
\item คงที่
\item เพิ่มเป็น 2 เท่า
\end{enumerate}
\textbf{คำตอบ: } ข
\textit{เฉลย: } N = mg cosθ เมื่อ θ เพิ่มจาก 30° เป็น 60° → cos60° < cos30° ดังนั้น N ลดลง
\HRule
\item[28] รถมวล 1000 kg แล่นด้วยความเร็ว 72 km/h บนถนนราบ ถูกเบรกจนหยุด ถ้า μₖ = 0.5 ระยะเบรกเป็นเท่าไหร่? (g = 10 m/s²)
\textbf{คำตอบ: } 40 m
\textit{วิธีทำ: } v = 72 km/h = 20 m/s
fk = μkN = 0.5 × 1000 × 10 = 5000 N
a = -fk/m = -5000/1000 = -5 m/s²
v² = u² + 2as → 0 = 400 + 2(-5)s → s = 400/10 = 40 m
\HRule
\item[29] เหตุใดจึงเปิดลิ้นชักได้ง่ายกว่าดันให้เคลื่อนที่ตอนเริ่มต้น?
\begin{enumerate}
\item เพราะน้ำหนักลิ้นชักลดลง
\item เพราะเสียดทานจลน์น้อยกว่าเสียดทานสถิต
\item เพราะแรงโน้มถ่วงลดลง
\item เพราะลิ้นชักเบาลง
\end{enumerate}
\textbf{คำตอบ: } ข
\textit{เฉลย: } เสียดทานสถิตสูงสุดมากกว่าเสียดทานจลน์ (μₛ > μₖ) ดังนั้นต้องออกแรงมากกว่าเพื่อเริ่มเคลื่อนที่ เมื่อเริ่มเคลื่อนแล้วเหลือแค่เสียดทานจลน์ที่น้อยกว่า จึงเปิดได้ง่ายกว่า
\HRule
\item[30] วัตถุมวล 10 kg วางบนพื้นเอียง 30° ต้องออกแรงขนานพื้นเอียงเท่าไหร่เพื่อดึงให้เคลื่อนที่ขึ้นด้วยความเร็วคงที่? (μₖ = 0.2, g = 10 m/s²)
\textbf{คำตอบ: } 56.9 N
\textit{วิธีทำ: } เมื่อเคลื่อนที่ขึ้นด้วยความเร็วคงที่: F = mg sinθ + fₖ
N = mg cosθ = 10×10×(√3/2) = 86.6 N
fₖ = μₖN = 0.2×86.6 = 17.32 N
F = 10×10×0.5 + 17.32 = 50 + 17.32 = 67.32 N... หรือปรับค่าตัวเลข: F = mg sinθ + μₖ mg cosθ = 10×10(sin30° + 0.2×cos30°) = 100(0.5 + 0.2×0.866) = 100(0.5+0.173) = 67.3 N
\HRule
\end{enumerate}
\end{document}
